% TOP_PROBLEM: {"id": 5300039, "title": "Infinite intersections of small Mandelbrot sets", "category_id": 11, "category": {"id": 11, "name": "dynamical_systems", "display_name": "Dynamical Systems", "description": "Problems about long-term behavior of deterministic systems, Hamiltonian dynamics, and periodic orbits.", "slug": "dynamical-systems", "order_index": 11, "created_at": "2026-07-31T15:26:25.671Z"}, "status": "partially_solved"}
% FIRST_POSTED: null
% ENTRY_KIND: statement_only
% Imported from: batch05.tex; UnsolvedMath ID 5300039
\documentclass[11pt]{article}
\usepackage[a4paper,margin=25mm]{geometry}
\usepackage{amsmath,amssymb,mathtools}
\usepackage{fontspec,unicode-math}
\setmainfont{Latin Modern Roman}
\setsansfont{Latin Modern Sans}
\setmathfont{Latin Modern Math}
\usepackage{xurl,hyperref}
\hypersetup{colorlinks=true,linkcolor=blue,urlcolor=blue,pdftitle={UnsolvedMath Problem 5300039}}
\newcommand{\sref}[2]{\hyperlink{source:#1:#2}{[S#2]}}
\newcommand{\cataloguescope}[1]{\paragraph{Mathematical question.}#1}
\begin{document}
% UnsolvedMath ID 5300039; source catalogue code AMR-052-0039
\setcounter{section}{5300038}
\section{Infinite intersections of small Mandelbrot sets}\label{5300039}
{\small\sffamily UnsolvedMath ID 5300039 · Dynamical Systems}
\subsection{Definitions and mathematical statement}

\paragraph{Shared notation.}$\mathbb N=\{1,2,\ldots\}$, and $\mathbb Z,\mathbb Q,\mathbb R,\mathbb C$ denote the integers, rationals, reals and complex numbers. Any other conventions are specified locally.


For a polynomial $f$ of degree at least two, $f^n$ is its $n$-fold iterate, $K(f)=\{z\in\mathbb C:(f^n(z))_{n\geq0}\text{ is bounded}\}$ is its filled Julia set, and $J(f)=\partial K(f)$ is its Julia set. A space is locally connected if each point has a neighbourhood basis of connected sets. Let $M=\{c\in\mathbb C:0\in K(z^2+c)\}$ be the Mandelbrot set. A hyperbolic component is a connected component of the parameters having an attracting periodic orbit. For a component $H$ of period at least two, let $\tau_H:M\to M$ be its Douady--Hubbard tuning embedding, whose image is the small copy containing $H$; write $H*c=\tau_H(c)$. A quadratic is renormalizable if an iterate $f_c^p$, $p\geq2$, restricts to a degree-two polynomial-like map $U\to V$, where $U,V$ are topological disks, $\overline U\subset V$, the map is proper, and its filled set contains $0$. Infinitely renormalizable parameters admit arbitrarily deep nested such renormalizations. A set is totally disconnected if all its connected components are singletons. The Hausdorff dimension $\dim_H E$ is the threshold exponent for Hausdorff measures $\mathcal H^s(E)$, defined using covers by sets of small diameter and sums of their $s$th powers. Planar Lebesgue measure is denoted by $m_2$.
\paragraph{Statement recorded in the supplied source.}
Does every nested intersection $\bigcap_{k\geq1}\tau_{H_1}\circ\cdots\circ\tau_{H_k}(M)$ reduce to one point? In the equivalent formulation asked by the source, is the set of infinitely renormalizable parameters totally disconnected? Does this set have Lebesgue measure zero, and does it have small Hausdorff dimension? \sref{5300039}{2}
\subsection{Short English statement}
Do all infinite nests of tuned Mandelbrot copies shrink to one parameter? Are infinitely renormalizable parameters totally disconnected, of area zero, and of small Hausdorff dimension?
\subsection{Sources}
\begin{description}
\item[\hypertarget{source:5300039:1}{S1}] ULAM.AI and the UnsolvedMath Contributors, UnsolvedMath: A Curated Collection of Open Mathematics Problems, record 5300039 (AMR-052-0039). CC BY 4.0. \url{https://huggingface.co/datasets/ulamai/UnsolvedMath}.
\item[\hypertarget{source:5300039:2}{S2}] John Milnor, ``Problems on local connectivity'', in Ben Bielefeld and Mikhail Lyubich (eds.), Problems in Holomorphic Dynamics (1992), Problem 14, pp. 22--23; arXiv:math/9205209v1. \url{https://arxiv.org/pdf/math/9205209}.
\end{description}
\label{end:5300039}
\end{document}
